\documentclass[11pt]{article}

\usepackage[letterpaper,margin=1in]{geometry}
\usepackage{amsmath,amssymb}
\usepackage{booktabs}
\usepackage{longtable}
\usepackage{array}
\usepackage{xcolor}
\usepackage{listings}
\usepackage{url}
\usepackage{hyperref}
\usepackage{titlesec}
\usepackage{parskip}
\usepackage{fancyhdr}
\usepackage{enumitem}

\hypersetup{
  colorlinks=true,
  linkcolor=black,
  urlcolor=blue!50!black,
  citecolor=black,
  pdftitle={LonghornSilicon Documentation Library --- Proposal},
  pdfauthor={LonghornSilicon},
  pdfsubject={Documentation Library Proposal, Version library-0.1}
}

\titleformat{\section}{\normalfont\Large\bfseries}{\thesection.}{0.6em}{}
\titleformat{\subsection}{\normalfont\large\bfseries}{\thesubsection}{0.6em}{}
\titleformat{\subsubsection}{\normalfont\normalsize\bfseries}{\thesubsubsection}{0.6em}{}

\pagestyle{fancy}
\fancyhf{}
\lhead{\small LonghornSilicon Documentation Library}
\rhead{\small \texttt{library-0.1}}
\cfoot{\thepage}
\renewcommand{\headrulewidth}{0.4pt}

\setlist[itemize]{leftmargin=*,topsep=2pt,itemsep=1pt}

\title{\textbf{LonghornSilicon Documentation Library \\
Proposal and Inventory}}
\author{LonghornSilicon}
\date{Version \texttt{library-0.1} \quad 2026-05-22}

\begin{document}
\maketitle
\thispagestyle{fancy}

\begin{abstract}
\noindent
LonghornSilicon now spans nineteen repositories and ships at least
fourteen formal LaTeX documents --- five published or draft papers,
seven block-level specifications, and an emerging set of integration
findings. The volume justifies a single canonical index. This
document proposes folding the library role into the existing
\texttt{LonghornSilicon/architecture} repository (rather than creating
a new one), inventories every formal document across the
organization, and lays out a four-step implementation plan to land
a public-facing index within one work session.
\end{abstract}

\section{Motivation}

The organization's documentation footprint grew faster than its
indexing scheme. A new collaborator (university partner, foundry
contact, journal reviewer, future student) currently has to discover
documents repo by repo. We have neither:

\begin{itemize}
\item A reading order across blocks (which paper covers the chip,
which covers each block, which covers compression, which covers
integration).
\item A status register (which papers are submitted, which are drafts,
which are superseded).
\item A shared bibliography file so that all internal papers cite each
other consistently.
\item A public-facing entry point for an external visitor to grasp
the program's scope without cloning each repo.
\end{itemize}

The cost of each missing piece is small in isolation but compounding.
The Muse Semi meeting prep on 2026-05-21 is a recent example: the
relevant artifacts (Sky130 sign-off, paper, ISA, integration audit,
the TurboQuant+ draft) lived in five separate places, and the
``what should I show them'' question required reconstructing the map
from memory rather than reading it off a page.

\section{Inventory}
\label{sec:inventory}

The tables below cover every formal LaTeX-source-bearing document
across the LonghornSilicon GitHub organization as of 2026-05-22. Each
row shows the canonical source repository, the path within that
repository, current status, and the document's role.

\subsection{System-level / architecture}

\begin{center}
\renewcommand{\arraystretch}{1.2}
\setlength{\tabcolsep}{4pt}
\begin{tabular}{@{}p{0.30\textwidth}p{0.30\textwidth}lp{0.20\textwidth}@{}}
\toprule
\textbf{Document} & \textbf{Source} & \textbf{Status} & \textbf{Role} \\
\midrule
$\lambda$ chip architecture & \texttt{architecture/paper/lambda.tex}
  & draft  & full chip overview \\
Dataflow walkthrough        & \texttt{architecture/dataflow\_walkthrough.md}
  & living & block interaction explainer \\
Architecture status         & \texttt{architecture/STATUS.md}
  & living & per-block development state \\
Floorplan visualisation     & \texttt{architecture/floorplan.html}
  & living & visual aid \\
\bottomrule
\end{tabular}
\end{center}

\subsection{Block 1 --- ACU (Attention Compute Unit)}

\begin{center}
\renewcommand{\arraystretch}{1.2}
\setlength{\tabcolsep}{4pt}
\begin{tabular}{@{}p{0.30\textwidth}p{0.30\textwidth}lp{0.20\textwidth}@{}}
\toprule
\textbf{Document} & \textbf{Source} & \textbf{Status} & \textbf{Role} \\
\midrule
Adaptive Precision Attention & \texttt{attention-compute-unit/}
                                \texttt{paper/adaptive\_precision\_attention.tex}
  & draft  & main block paper \\
Precision Controller ISA     & \texttt{attention-compute-unit/}
                                \texttt{docs/isa/precision\_controller\_isa.tex}
  & v0.1 stable & interface specification \\
MAC array design             & \texttt{attention-compute-unit/}
                                \texttt{docs/mac\_array\_design.tex}
  & v0.1 stable & sub-block design \\
Reference model API          & \texttt{attention-compute-unit/}
                                \texttt{docs/reference\_model\_api.tex}
  & v0.1 stable & compiler-facing API \\
SW overview                  & \texttt{attention-compute-unit/}
                                \texttt{docs/sw\_overview.tex}
  & v0.1 stable & compiler-team orientation \\
Meeting handout              & \texttt{attention-compute-unit/}
                                \texttt{docs/meeting\_handout.tex}
  & event piece & one-pager (compiler meeting) \\
\bottomrule
\end{tabular}
\end{center}

\subsection{Block 2 --- KV Cache Engine}

\begin{center}
\renewcommand{\arraystretch}{1.2}
\setlength{\tabcolsep}{4pt}
\begin{tabular}{@{}p{0.30\textwidth}p{0.30\textwidth}lp{0.20\textwidth}@{}}
\toprule
\textbf{Document} & \textbf{Source} & \textbf{Status} & \textbf{Role} \\
\midrule
KV Cache Engine ISA       & \texttt{kv-cache-engine/docs/isa/}
                            \texttt{kv\_cache\_engine\_isa.tex}
  & v0.1 stable & interface specification \\
Reference model API       & \texttt{kv-cache-engine/docs/}
                            \texttt{reference\_model\_api.tex}
  & v0.1 stable & compiler-facing API \\
SW overview               & \texttt{kv-cache-engine/docs/}
                            \texttt{sw\_overview.tex}
  & v0.1 stable & compiler-team orientation \\
TurboQuant+ paper         & \texttt{turboquant/} (in progress)
  & draft, private & compression algorithm paper \\
\bottomrule
\end{tabular}
\end{center}

\subsection{Cross-block findings}

\begin{center}
\renewcommand{\arraystretch}{1.2}
\setlength{\tabcolsep}{4pt}
\begin{tabular}{@{}p{0.30\textwidth}p{0.30\textwidth}lp{0.20\textwidth}@{}}
\toprule
\textbf{Document} & \textbf{Source} & \textbf{Status} & \textbf{Role} \\
\midrule
KVCE $\times$ ACU integration audit & \texttt{attention-compute-unit/}
                                  \texttt{docs/findings/}
                                  \texttt{kvce\_acu\_integration\_audit.tex}
  & v0.2 current & integration test + remediation \\
Cross-block conflict register      & \texttt{attention-compute-unit/}
                                  \texttt{docs/findings/}
                                  \texttt{kvce\_acu\_architectural\_conflicts.tex}
  & v0.1 (forthcoming) & standing design-tension register \\
Phase~1 entropy finding            & \texttt{attention-compute-unit/}
                                  \texttt{docs/findings/}
                                  \texttt{phase1-entropy-finding.md}
  & frozen & evolutionary policy discovery \\
Sparsity controller finding        & \texttt{attention-compute-unit/}
                                  \texttt{docs/findings/}
                                  \texttt{sparsity-controller-finding.md}
  & shelved & XAttention proxy negative result \\
\bottomrule
\end{tabular}
\end{center}

\subsection{Adjacent research}

\begin{center}
\renewcommand{\arraystretch}{1.2}
\setlength{\tabcolsep}{4pt}
\begin{tabular}{@{}p{0.30\textwidth}p{0.30\textwidth}lp{0.20\textwidth}@{}}
\toprule
\textbf{Document} & \textbf{Source} & \textbf{Status} & \textbf{Role} \\
\midrule
Don't Waste Bits FPGA controller & \texttt{dont-waste-bits/research/}
                                    \texttt{paper/fpga\_controller\_paper.tex}
  & draft, submission-ready & variable-precision controller paper \\
DWB original (verified)          & \texttt{dont-waste-bits/2604.04722v1.pdf}
  & external (CVPR~2026) & reference \\
LASSO KV coprocessor             & \texttt{lasso/writing\_outputs/}
                                    \texttt{lasso\_paper.tex}
  & draft & Sky130 KV-cache compressor \\
HADES step~5                     & \texttt{decode-pipeline-optimizer/}
                                    \texttt{paper/hades\_step5.tex}
  & draft & speculative decoding assist \\
FP4 multiplier                   & \texttt{fp4-multiplier/paper.pdf}
  & PDF only & minimum-gate FP4 (E2M1) \\
Chipathon microarchitecture      & \texttt{Chipathon/docs/}
                                    \texttt{Microarchitecture.pdf}
  & PDF only & SSCS Chipathon 130nm design \\
\bottomrule
\end{tabular}
\end{center}

\subsection{Tooling / adjacent}

Repositories with no LaTeX documents but supporting the library's
context: \texttt{architecture/.github/}, \texttt{website},
\texttt{.github}, \texttt{SRAM\_16384x32}, \texttt{kelle-simulator},
\texttt{kelle-fpga}, \texttt{kv-accelerator-comparison},
\texttt{OpenLane}, \texttt{openlane2}, \texttt{skywater-pdk},
\texttt{tiny-gpu}. These appear in the library only as cross-reference
targets, not as document sources.

\section{Three structural options}

\subsection{Option A: dedicated \texttt{LonghornSilicon/library} repository}

Create a new repository whose sole purpose is the canonical index.
Contains a master \texttt{index.tex}, a shared \texttt{references.bib},
and a GitHub Pages render.

\begin{description}[leftmargin=2em,style=nextline]
\item[Pros]
Clean separation of concerns; library has no other purpose so it
cannot drift into being half-architecture, half-index. Discoverable
as ``the library'' from the organization landing page.
\item[Cons]
Yet another repository to maintain. Risk of becoming a duplicate
storage of papers rather than a pure index. Cross-references from
existing repos must be added by hand.
\end{description}

\subsection{Option B: extend the existing \texttt{architecture} repository}

\texttt{LonghornSilicon/architecture} already exists, already holds
the system-level $\lambda$ paper, the dataflow walkthrough, the
status file, and the floorplan visualisation. Add to it a
\texttt{papers/} directory with the library index and shared
bibliography.

\begin{description}[leftmargin=2em,style=nextline]
\item[Pros]
Zero new repositories. The architecture repo's existing role
(cross-block reference) generalises naturally to ``cross-block
documentation index''. The $\lambda$ paper already pulls together
all block claims --- the library is its natural sibling. GitHub
Pages can render directly from this repository.
\item[Cons]
The architecture repo's scope widens; ``architecture'' now also
means ``publication index''. Mitigated by clear directory boundaries
(\texttt{paper/} for $\lambda$, \texttt{papers/} for the library
index, \texttt{specs/} for cross-block specs).
\end{description}

\subsection{Option C: distributed cross-references, no new home}

Add a \texttt{LIBRARY.md} or \texttt{INDEX.md} to each repository
that cross-links to the others. No central index.

\begin{description}[leftmargin=2em,style=nextline]
\item[Pros]
Lowest effort. Each repo stays self-contained.
\item[Cons]
No single entry point for an external visitor. Cross-references must
be kept in sync across nineteen repositories. Does not address the
``which paper covers what'' question for someone scanning from outside.
\end{description}

\subsection{Recommendation}

\textbf{Option B}: fold the library into
\texttt{LonghornSilicon/architecture}. The architecture repo is
already the de-facto cross-cutting home; the $\lambda$ paper is
already there; the dataflow walkthrough is already there; adding the
index closes the loop without creating new repositories. Option~A
is the second-best choice if a clean separation of concerns
becomes preferred later; Option~C is rejected as not solving the
external-visitor problem.

\section{Proposed structure within \texttt{architecture}}

\begin{verbatim}
LonghornSilicon/architecture/
|-- README.md                  # entry point, links to papers/INDEX
|-- STATUS.md                  # existing per-block development status
|-- arch.yml                   # existing structured architecture model
|-- dataflow_walkthrough.md    # existing dataflow explainer
|-- floorplan.html             # existing floorplan visualisation
|-- paper/
|   `-- lambda.tex             # existing system-level paper
|-- papers/                    # NEW: library
|   |-- INDEX.md               # human-readable index, rendered to Pages
|   |-- index.tex              # LaTeX-typeset companion (this proposal's
|   |                          #   document table, kept in sync)
|   |-- references.bib         # NEW: shared bibliography for all
|   |                          #   LonghornSilicon papers to cite
|   `-- abstracts/
|       |-- adaptive_precision_attention.md
|       |-- precision_controller_isa.md
|       |-- kv_cache_engine_isa.md
|       |-- turboquant.md
|       |-- fpga_controller_paper.md
|       |-- lasso_paper.md
|       |-- hades_step5.md
|       |-- kvce_acu_integration_audit.md
|       `-- ...
|-- specs/                     # NEW: convenience links to canonical
|   |                          #   block specs (symlinks or short files)
|   `-- README.md              # ``each block's ISA + ref model API''
`-- findings/                  # NEW: cross-block findings sourced from
    |                          #   block repos
    `-- README.md              # links to integration audit etc.
\end{verbatim}

The \texttt{papers/} directory is the library proper. Each entry
in \texttt{papers/abstracts/} is a one-page markdown summary
(canonical title, authors, status, abstract, link to the
authoritative source repo). The \texttt{INDEX.md} aggregates these
into a single rendered page; \texttt{index.tex} is its LaTeX-typeset
form for inclusion in funding reports or partnership packages.

\section{Shared bibliography}

The single highest-leverage artifact is
\texttt{papers/references.bib}. Today, when the
\texttt{adaptive\_precision\_attention} paper wants to cite
TurboQuant+, the author hand-writes a BibTeX stanza. When the
KVCE \texttt{lasso\_paper} wants to cite the precision controller,
another hand-written stanza appears, potentially with a different
author list or title formatting. A canonical
\texttt{references.bib} with stanzas like:

\begin{lstlisting}[basicstyle=\ttfamily\small,frame=single]
@misc{lhs_acu_precision_controller,
  title  = {Adaptive Precision Attention: Evolutionary
            Discovery to Hardware-Verified Ratio Gates},
  author = {Talasila, Chaithu and {LonghornSilicon}},
  year   = {2026},
  url    = {https://github.com/LonghornSilicon/
            attention-compute-unit},
  note   = {Block 1 of the LonghornSilicon LLM
            inference accelerator},
}

@misc{lhs_kvce_turboquant,
  title  = {TurboQuant+: Asymmetric K/V Compression
            for LLM KV Caches},
  ...
}
\end{lstlisting}

allows every paper in the organization to cite every other paper
consistently. Future-us, advisors, foundry partners, and reviewers
all encounter the same author lists and titles for the same work.

\section{Implementation plan}

Four steps, each landable in under thirty minutes by a single
session.

\begin{enumerate}
\item \textbf{Bootstrap the \texttt{papers/} directory in
\texttt{architecture}}. Add the directory, the \texttt{INDEX.md}
seeded from \S\ref{sec:inventory}, the \texttt{abstracts/}
sub-directory with one stub abstract per document, and the empty
\texttt{references.bib}.

\item \textbf{Populate abstracts.} For each document in
\S\ref{sec:inventory}, copy the existing abstract (or write one
where missing) into the corresponding
\texttt{papers/abstracts/<name>.md} file. Link to the canonical
source.

\item \textbf{Populate the shared bibliography.} Add one BibTeX
stanza per document to \texttt{references.bib}, with consistent
author, title, and URL fields. Update existing papers to
\verb|\bibliography{../architecture/papers/references}| (or copy
the bib into each paper as needed).

\item \textbf{Render to GitHub Pages.} Enable Pages on
\texttt{architecture} pointing at \texttt{papers/INDEX.md}. The
organization landing page links to it. Done.
\end{enumerate}

Total work to first usable state: approximately two hours.

\section{Open questions}

\begin{enumerate}
\item Should TurboQuant+ live in \texttt{turboquant} (current home)
or be promoted to a dedicated repository (e.g.\ \texttt{turboquant-plus})
now that it is the canonical KVCE compression algorithm?
\item Should the library index also enumerate \emph{external}
references the LonghornSilicon papers consistently cite (FlashAttention
1/2/3, GEAR, KIVI, KVQuant, RotateKV, SHIELD, etc.) so that
the shared \texttt{references.bib} stops requiring hand-additions
in each paper?
\item Does the website (\texttt{LonghornSilicon/website}) consume the
library directly (e.g.\ render the index in-place) or maintain its
own publication list that pulls from the library?
\item Status semantics: do we adopt three states
(\textit{draft}, \textit{submitted}, \textit{accepted}) or finer
granularity (\textit{draft / under-review / accepted / published /
superseded / archived})?
\item Versioning: do papers carry their own version codes
(\texttt{pc-isa-0.1}, \texttt{kvce-isa-0.1}, \texttt{kvce-acu-audit-0.2})
into the library entries, or does the library normalise to a single
versioning scheme?
\end{enumerate}

These are not blockers; they can be resolved during step 2 of the
implementation plan.

\section*{Change log}

\begin{itemize}
\item \texttt{library-0.1} (2026-05-22): inventory completed
across nineteen organization repositories; structural options
evaluated; recommendation made; implementation plan drafted.
\end{itemize}

\end{document}
